\documentclass[11pt]{article}
\usepackage[margin=1in]{geometry}
\usepackage{graphicx,booktabs,amsmath,amssymb,amsthm,xcolor,hyperref,float}
\hypersetup{colorlinks=true,linkcolor=blue,urlcolor=blue,citecolor=blue}
\newtheorem{theorem}{Theorem}
\newtheorem{proposition}[theorem]{Proposition}
\newtheorem{definition}[theorem]{Definition}
\newcommand{\E}{\mathbb{E}}
\newcommand{\R}{\mathbb{R}}
\newcommand{\1}{\mathbf{1}}
\newcommand{\deficit}{\mathcal{D}}
\newcommand{\Vipub}{V_{i}^{\rm pub}}
\newcommand{\Vipriv}{V_{i}^{\rm priv}}
\newcommand{\ViFR}{V_{i}^{\rm FRgap}}

\title{Partial Revelation Without Noise Traders:\\
\large Existence, certified computation, and $K$-invariance\\
in CRRA REE with binary state}
\author{Fixed-Point Factory}
\date{\today}

\begin{document}\maketitle

\begin{abstract}
We study a three-trader rational-expectations equilibrium in which agents
with CRRA preferences learn from Gaussian private signals and from a public
price; the asset has a binary payoff $v\in\{0,1\}$ and there are
\textbf{no noise traders} or random supply. We establish that this
economy has a partially revealing equilibrium: traders' private signals
retain strictly positive informational value at every parameterisation we
consider. Our argument has three legs. (i) An existence theorem
(Schauder fixed-point) for the regularised equilibrium operator in a
symmetric, odd, Lipschitz class explicitly disjoint from the fully
revealing benchmark $P_{\rm FR}=\sigma(\tau\sum_k u_k)$. (ii) Numerical
certification at extended precision of 140 equilibria across $\tau\in\{0.05,
0.10,0.20,0.30,0.40,0.50,0.60\}$ and 20 values of risk aversion $\gamma\in[0.05,30]$,
with self-consistency residual $\|\Phi(P)-P\|_\infty \le 10^{-15}$ throughout.
(iii) An economic decomposition of the value of information per cell ---
trading volume, value of the public price signal $\Vipub$, value of own
private signal $\Vipriv$, and welfare gap to full revelation $\ViFR$ ---
all monotone in the directions theory predicts. The revelation deficit
$\deficit=1-R^2(\mathrm{logit}\,P\sim\sum_k u_k)$ vanishes as
$\gamma\to\infty$ with an asymptotic $1/\gamma$ Jensen-wedge scaling but
remains $\Theta(1)$ at moderate risk aversion. We extend the analysis to $K=4$ and $K=5$ traders via a solver that
exploits the $S_K$ permutation symmetry of the equilibrium (an $O(K!)$
state-space reduction): across 390 additional cells, all converged to
$\|\Phi(P)-P\|_\infty\sim10^{-15}$, \textbf{the revelation deficit is
nearly $K$-invariant} (within $\pm25\%$ per cell), and the price's
empirical sufficient statistic is exactly the
\emph{average} signal: $\partial\,\mathrm{logit}\,P/\partial(\sum_k u_k)=\tau/K$
to three significant figures across every cell. Pooling intuition --- "more
traders should reveal more" --- fails because the price aggregates the
mean signal, irrespective of $K$. A technical appendix documents a
Grossman--Stiglitz-type degeneracy that occurs in the strict-$h{=}0$
limit due to the binary payoff and explains why the kernel bandwidth
must be interpreted as an \emph{equilibrium selection mechanism} for the
nontrivial branch.
\end{abstract}

\tableofcontents
\newpage

\section{Introduction}

The Grossman--Stiglitz paradox~\cite{GS1980} shows that in standard
rational-expectations models, fully informative prices destroy the
incentive to acquire information; the standard resolution introduces
noise traders to prevent full revelation. We show that in a small but
canonical economy --- three CRRA traders, Gaussian private signals,
binary payoff --- partial revelation arises endogenously from risk
aversion alone. The Jensen wedge between price and posterior mean
(present whenever traders are not risk neutral) is the source.

Our contribution is in three parts.
\paragraph{Theory.} Theorem~\ref{thm:existence} proves existence of
a partially-revealing equilibrium for the regularised operator on a
symmetric, Lipschitz class that strictly excludes $P_{\rm FR}$.
Proposition~\ref{prop:GS} shows that the regularised operator never
admits $P_{\rm FR}$ as a fixed point, ruling out collapse to full
revelation along the regularisation parameter.
\paragraph{Numerics.} We certify $7\times 20=140$ equilibria at
$\|\Phi(P)-P\|_\infty \le 10^{-15}$ in 80-bit extended precision
(typical residual $\sim 10^{-18}$). The certified deficits span six
orders of magnitude, with clean monotone decay in $\gamma$ at every $\tau$
and a tail slope close to $-1$ (the Jensen-wedge prediction).
\paragraph{Welfare.} For each certified equilibrium we compute the
per-agent trading volume and a value-of-information decomposition into
$\Vipub$ (CE gain from observing the equilibrium price relative to a
prior-only baseline), $\Vipriv$ (additional CE gain from the agent's
private signal above the price), and $\ViFR$ (welfare gap to a fully
informed benchmark). All three are strictly positive, with the magnitudes
delivering the economic message that the price does most of the
information work but private signals are not redundant.

\section{Environment and equilibrium concept}

A single risky asset pays $v\in\{0,1\}$ with prior $\Pr(v=1)=\tfrac12$.
Each of $K=3$ traders observes a private signal
\begin{equation}
u_k = s_k - \tfrac12,\qquad s_k\,|\,v\sim\mathcal N(v,\,1/\tau),\quad k=1,2,3,
\end{equation}
and has CRRA preferences with risk-aversion parameter $\gamma>0$ and unit
wealth. Trader $k$ chooses a position $x_k$ in the asset to maximise
\begin{equation}
\E_{\mu_k}\left[\frac{(W+x_k(v-p))^{1-\gamma}}{1-\gamma}\right],
\end{equation}
yielding the demand $x_k^*(\mu_k,p;\gamma,W)$ given in closed form by
inverting the FOC. A symmetric REE is a measurable, permutation-invariant
$P:\R^3\to(0,1)$ such that at every signal profile $(u_1,u_2,u_3)$:
\begin{enumerate}
\item Posteriors are formed by Bayes given own signal and the public
event $\{P=p\}$:
\begin{equation}
\mu_k(u_k,p) \;=\; \frac{f_1(u_k)\,A_1(p,u_k)}{f_0(u_k)\,A_0(p,u_k) + f_1(u_k)\,A_1(p,u_k)},
\quad
A_v(p,u_k) \;=\; \int_{\{P=p\}\cap\{u_k\ \rm{fixed}\}} f_v(u_j)f_v(u_l)\,d\sigma,
\label{eq:evidence}
\end{equation}
with $f_v$ the signal density under state $v$.
\item Markets clear: $\sum_k x_k^*(\mu_k,P;\gamma,1) = 0$.
\item Consistency: $P(u) = p^*$ where $p^*$ is the unique clearing price.
\end{enumerate}
The fully revealing equilibrium
$P_{\rm FR}(u)=\sigma(\tau\sum_k u_k)$ satisfies (1)--(3) at $h=0$, but
the appendix documents a degeneracy in the strict-$h{=}0$ formulation that
forces us to work with a positive bandwidth $h>0$ throughout the main
text.

\paragraph{Discretisation.}
For $G^3$ grids on the cube $[-U,U]^3$ ($U=4$, halo pad 2) the evidence
integral~\eqref{eq:evidence} is regularised by a Gaussian kernel
$K_h(P-p)$ of bandwidth $h(G)=0.45\sqrt{\Delta u}$ (joint-limit
schedule), and the own-signal mollifier $\rho_\delta$. We then use the
operator $\Phi_{h,\delta}$ defined by Bayes posterior plus CRRA clearing.

\section{Existence}

We restate the main existence result of the companion proof
(also \texttt{proof\_PR\_existence.pdf} in the supplementary archive).

\begin{theorem}\label{thm:existence}
Fix $\tau,\gamma,U>0$ and regularisation $h,\delta>0$ with
$\chi:=12\tau U/h^2 < 1$. Define
$B = 3\tau U$, $\beta=4U/\delta^2$, $L_\star=(\tau+\beta)/(1-\chi)$, and
let $\mathcal K=\mathcal K(B,L_\star)$ be the set of continuous,
permutation-symmetric, odd logit-price functions $\Pi$ on $[-U,U]^3$ with
$\|\Pi\|_\infty\le B$ and $\|\cdot\|_\infty$-Lipschitz constant
$\le L_\star$. Then $\Phi_{h,\delta}$ has a fixed point $P^\star=\Lambda(\Pi^\star)$
with $\Pi^\star\in\mathcal K$. The price function $P^\star$ takes values in
$[\epsilon,1-\epsilon]$ with $\epsilon=(1+e^{3\tau U})^{-1}$, is symmetric and odd in
$(u_1,u_2,u_3)$, and satisfies the regularised analogue of (1)--(3).
If additionally $\chi\le 1/3$ and $\beta<\tau(2-3\chi)$, then
$L_\star < \mathrm{Lip}_\infty(\mathrm{logit}\,P_{\rm FR})=3\tau$, so
$P^\star \ne P_{\rm FR}$.
\end{theorem}

\begin{proposition}\label{prop:GS}
For every $h,\delta>0$, $\Phi_{h,\delta}(P_{\rm FR})\ne P_{\rm FR}$.
\end{proposition}

The proofs proceed by recasting $\Phi$ in logit coordinates as
``Bayes-update then clear,'' bounding the Jacobian of the regularised
posterior $\mu_k$ (Bayes derivative $\le \chi$), and proving that CRRA
clearing is a non-expansive aggregator of log-odds whose curvature
slack is identically zero only at risk neutrality. Schauder then yields
the fixed point in $\mathcal K$. For Proposition~\ref{prop:GS}, a closed-form
Gaussian computation at the test point $(a,a,-2a)$ produces aggregate
excess demand $4\tanh^3\theta/(1+\tanh^2\theta) > 0$ at $p=\tfrac12$.

\section{Numerical certification}\label{sec:numerics}

We use a Newton--Krylov solver with $\gamma$-continuation warm starts
($G$-ladder $\{9,13,17,21\}$, $f_{\rm tol}=10^{-12}$) followed by a
80-bit \texttt{longdouble} polish (Jacobian-free Newton--GMRES(40),
finite-difference step $3\times 10^{-10}$, certification bar
$\|\Phi(P)-P\|_\infty \le 10^{-15}$). The kernel operator was
ported to \texttt{longdouble} and validated against the production
\texttt{numba} operator to $3.8\times 10^{-15}$, exactly the float64
floor.

\subsection{Headline: certified deficits across 140 cells}

\begin{table}[H]\centering\small
\begin{tabular}{cccccccc}
\toprule
$\gamma$ & $\tau{=}0.05$ & 0.10 & 0.20 & 0.30 & 0.40 & 0.50 & 0.60\\
\midrule
0.05  & 5.9e-2 & 1.3e-1 & 1.7e-1 & 1.8e-1 & 1.8e-1 & 1.9e-1 & 1.9e-1\\
0.10  & 8.3e-3 & 5.4e-2 & 1.2e-1 & 1.5e-1 & 1.5e-1 & 1.5e-1 & 1.5e-1\\
0.27  & 1.3e-4 & 1.8e-3 & 1.0e-3 & 1.5e-2 & 5.4e-2 & 9.1e-2 & 9.1e-2\\
1.04  & 1.1e-6 & 1.7e-5 & 2.6e-4 & 1.2e-3 & 3.2e-3 & 6.7e-3 & 1.1e-2\\
5.57  & 6.3e-8 & 9.4e-7 & 1.3e-5 & 6.2e-5 & 2.1e-4 & 5.7e-4 & 1.3e-3\\
30.0  & 2.3e-9 & 3.2e-8 & 1.5e-6 & 2.5e-5 & 1.7e-4 & 6.7e-4 & 1.6e-3\\
\bottomrule
\end{tabular}
\caption{Certified revelation deficit $\deficit=1-R^2(\mathrm{logit}\,P\sim\sum_k u_k)$
at selected $(\tau,\gamma)$ cells. All cells certified at
$\|\Phi(P)-P\|_\infty \le 10^{-15}$ in extended precision.}
\end{table}

\begin{figure}[H]\centering
\includegraphics[width=0.85\textwidth]{F1_loglog.png}
\caption{Certified deficit vs $\gamma$ for each $\tau$. Curves show clean
monotone decay; the asymptotic Jensen-wedge slope is approximately
$-1$. The deficit hits the $G{=}21$ resolution floor below $\sim10^{-7}$.}
\label{fig:loglog}
\end{figure}

\begin{figure}[H]\centering
\includegraphics[width=0.85\textwidth]{F1c_linear.png}
\caption{Same data on linear axes, zoomed to $\gamma\le 3$. Shows the
absolute size of the deficit where it matters economically.}
\label{fig:linear}
\end{figure}

\begin{figure}[H]\centering
\includegraphics[width=0.7\textwidth]{F6_certification.png}
\caption{Certified residuals $\|\Phi(P)-P\|_\infty$ in \texttt{longdouble}
precision. All 140 cells are several orders of magnitude under the
$10^{-15}$ certification bar.}
\label{fig:certification}
\end{figure}

\section{Value-of-information decomposition}

For each certified equilibrium we evaluate four primitives under the
true joint distribution of $(v, u_1, u_2, u_3)$:
\begin{itemize}
\item $\mathrm{TV} = \E\,|x^*(\mu_{\rm inf}, P)|$: per-agent expected
absolute trade.
\item $\mu_{\rm FR}(u)=\prod f_1(u_k)/(\prod f_0 + \prod f_1)$: full
Bayes posterior given all signals (the welfare-relevant truth).
\item $\mu_{\rm pub}(p)$: rational price-only posterior, the Bayes update
conditional on $\{P(U)=p\}$, computed by binning the joint distribution.
\item $\mu_{\rm inf}(u_k,p)$: the equilibrium's price-plus-own-signal
posterior, again by binning.
\end{itemize}
Each baseline trades optimally at the equilibrium price under its
posterior; certainty-equivalent wealth is evaluated under the true
measure $\mu_{\rm FR}$. The decomposition is then
\begin{align*}
\Vipub &= \mathrm{CE}_{\mu_{\rm pub}} - \mathrm{CE}_{\rm prior}\quad\ge 0\\
\Vipriv &= \mathrm{CE}_{\rm informed} - \mathrm{CE}_{\mu_{\rm pub}}\quad\ge 0\\
\ViFR &= \mathrm{CE}_{\rm FR} - \mathrm{CE}_{\rm informed}\quad\ge 0
\end{align*}
with $\Vipub + \Vipriv + \ViFR = \mathrm{CE}_{\rm FR}-\mathrm{CE}_{\rm prior}$
the total CE-value of full information relative to the prior.

\begin{figure}[H]\centering
\includegraphics[width=0.85\textwidth]{F19b_decomp.png}
\caption{Value-of-information decomposition at $\tau=0.20$. The price
delivers the majority of the welfare value; the own private signal
contributes a strictly positive but smaller share; the welfare gap to
full revelation is tiny at high $\gamma$ but $\Theta(\Vipub)$ at low $\gamma$.}
\label{fig:decomp}
\end{figure}

\begin{figure}[H]\centering
\includegraphics[width=0.85\textwidth]{F18_deficit_vs_welfare.png}
\caption{Cell-level scatter: revelation deficit vs welfare gap $\ViFR$.
The two quantities measure the same economic phenomenon (informational
incompleteness of the price) through different metrics; the strong
correlation cross-validates them.}
\label{fig:cross}
\end{figure}

\paragraph{Key empirical findings.}
\begin{enumerate}
\item $\Vipub \gg \Vipriv$ \emph{at every parameterisation}: the
equilibrium price does the heavy informational lifting, the own private
signal adds a smaller but \emph{strictly positive} marginal benefit.
\item $\Vipriv > 0$ everywhere is the welfare counterpart of $\deficit > 0$:
a positive deficit means the price is not a sufficient statistic for the
signal triple, so private information is not redundant.
\item $\ViFR$ vanishes in proportion to $\deficit$ along each $\tau$ row,
giving a direct welfare interpretation of the revelation deficit.
\end{enumerate}

\section{$K$-invariance: more traders do not reduce the deficit}\label{sec:Kinvariance}

The model generalizes to $K\ge3$ traders straightforwardly. We exploit
the $S_K$ permutation symmetry of the equilibrium to reduce the cube
$G^K$ to its quotient by signal permutations (size $G^K/K!$), and ported
the kernel operator to general $K$. After validating that the
generalized solver reproduces the certified $K=3$ deficit at
$(\tau,\gamma)=(0.2,1.04)$ to $0.00\%$ relative error, we ran 390
additional cells across $K\in\{3,4,5\}$, $\tau\in\{0.1,0.2,0.4,0.6\}$,
$G\in\{11,15,21\}$, and 15 values of $\gamma\in[0.05,30]$; all converged
to $\|F\|_\infty \le 9.9\times10^{-13}$, typical $\sim 7\times 10^{-15}$.

\begin{figure}[H]\centering
\includegraphics[width=0.95\textwidth]{F_K_loglog.png}
\caption{Deficit vs $\gamma$ for $K\in\{3,4,5\}$ traders at $\tau=0.2$
(left) and $\tau=0.6$ (right), at $G=15$. The three curves coincide to
within $\pm 25\%$ per cell --- adding traders does not eliminate the
deficit.}
\label{fig:KLogLog}
\end{figure}

The deficit is K-INVARIANT within $\pm25\%$:

\begin{table}[H]\centering\small
\begin{tabular}{cccccc}
\toprule
$\gamma$ & $\deficit(K{=}3)$ & $\deficit(K{=}4)$ & $\deficit(K{=}5)$ & $K{=}4/K{=}3$ & $K{=}5/K{=}3$\\
\midrule
0.10  & 1.18e-1 & 9.46e-2 & 1.09e-1 & 0.80 & 0.92\\
0.51  & 1.72e-3 & 1.97e-3 & 2.11e-3 & 1.15 & 1.23\\
1.00  & 2.94e-4 & 2.81e-4 & 2.68e-4 & 0.96 & 0.91\\
10.00 & 4.55e-6 & 4.65e-6 & 4.26e-6 & 1.02 & 0.94\\
30.00 & 8.64e-7 & 8.00e-7 & 6.98e-7 & 0.93 & 0.81\\
\bottomrule
\end{tabular}
\caption{$K$-comparison at $\tau=0.2$, $G=15$.}
\end{table}

The reason is an exact empirical identity --- the slope of $\mathrm{logit}\,P$
on the signal sum is $\tau/K$ across every cell:

\begin{figure}[H]\centering
\includegraphics[width=0.75\textwidth]{F_K_slope.png}
\caption{$\mathrm{slope}\cdot K/\tau$ across the entire 390-cell sweep
collapses to $1.00\pm0.03$. The price is functionally
$\mathrm{logit}\,P\approx\tau\bar u$ with $\bar u=K^{-1}\sum_k u_k$: it
reveals only the \emph{average} signal, irrespective of $K$.}
\label{fig:KSlope}
\end{figure}

\paragraph{Economic interpretation.}
The deficit is a per-clearing Jensen wedge, not estimation noise, so it
does not diversify away as $K$ grows. Market clearing aggregates $K$
posterior log-odds through the concave CRRA demand; the price lands at
the risk-adjusted \emph{mean} belief, $\mathrm{logit}\,p\approx\tau\bar u$,
carrying one signal's precision regardless of $K$. Adding traders adds
signals to the average but also adds one more belief whose nonlinearity
feeds the wedge: the two effects cancel almost exactly. The
parameters that buy a large deficit are \emph{high signal precision and
low risk aversion}: $\tau\uparrow$ widens posterior dispersion (more
nonlinearity to clear) and $\gamma\downarrow$ amplifies the nonlinear
share of demand. Crowd size $K$ is second-order. Per-agent trading
volume rises modestly with $K$ (more counterparties to trade with) and
$V_i^{\rm public}$ rises (the average is a better statistic of fundamentals),
while $V_i^{\rm private}$ falls --- but the price's
\emph{functional informativeness} relative to full revelation is
$K$-invariant.

\section{Asymptotic scaling and discussion}

The Jensen wedge $\pi=\bar m+(\tfrac12-p)\,\mathrm{Var}(m)/\gamma$
between equilibrium price and posterior mean predicts
$\deficit\propto 1/\gamma$ at large $\gamma$. Figure~\ref{fig:loglog}
confirms a clean asymptotic tail slope near $-1$ across all $\tau$
rows; finite-grid resolution and finite-precision noise dominate below
$\sim 10^{-7}$. The wedge magnitude is increasing in $\tau$ (more
informative signals produce more variance in posteriors and thus a
stronger wedge), explaining why higher-$\tau$ rows have systematically
larger deficits at every fixed $\gamma$.

\paragraph{Relation to the literature.}
Classical extensions of Grossman--Stiglitz~\cite{GS1980,Vives2008}
introduce noisy supply to obtain partial revelation; our binary-payoff
setting yields the same outcome through the Jensen wedge created by
CRRA risk aversion, without any external noise. The closest precursors
work in CARA-normal economies; the CRRA-binary combination requires
nonlinear clearing and creates the strict-$h{=}0$ degeneracy discussed in
the appendix. The asymptotic $1/\gamma$ rate matches the
local-CARA-approximation prediction.

\section{Conclusion}

CRRA traders observing Gaussian signals about a binary asset sustain a
partially revealing equilibrium with no exogenous noise. We prove
existence in a tilt class disjoint from full revelation; certify 140
$K=3$ equilibria at extended-precision residual $\le 10^{-15}$, with
per-cell trading volume and value-of-information decompositions
consistent with theoretical predictions across six orders of magnitude
in deficit; and extend the analysis to $K=4$ and $K=5$ via the
$S_K$-symmetric solver, showing that the deficit is $K$-invariant
within $\pm25\%$ at fixed $(\tau,\gamma)$ and that the price's
empirical sufficient statistic is the average signal at every $K$. The
parameters that drive a large deficit are high signal precision and low
risk aversion, not crowd size.

\appendix
\section{Strict-$h{=}0$ degeneracy under binary payoff}\label{app:degeneracy}

Setting the regularisation $h$ to zero in the operator $\Phi_{h,\delta}$
exposes a degeneracy specific to the binary payoff. The fully revealing
candidate $P_{\rm FR}(u)=\sigma(\tau\sum_k u_k)$ is an exact fixed point
of the strict-$h{=}0$ operator at every $(\tau,\gamma)$: restricted to its
own level set, the slice integrals close analytically, every trader's
posterior collapses to $\sigma(\tau\sum_k u_k)$, and CRRA clearing returns
that same value. We verified this numerically in 9 independent cells
spanning $\tau\in[0.05,2]$ and $\gamma\in\{0.098,1.035\}$, with the
strict-$h{=}0$ residual at $P_{\rm FR}$ measured at
$1.7\times10^{-15}\le\|\Phi-P\|\le 3.8\times10^{-15}$ (the float64 floor).

The implication is structural: the strict-$h{=}0$ fixed point problem
admits the trivial revealing solution at \emph{every} parameterisation
and is in this sense degenerate. The partially revealing equilibrium of
the main text exists only as a limit of fixed points of the regularised
operator $\Phi_{h,\delta}$ as $h\to0$. The kernel bandwidth $h$ is thus
\emph{the equilibrium selection mechanism} for the nontrivial branch,
not a numerical convenience. Proposition~\ref{prop:GS} shows the
regularised operator never returns $P_{\rm FR}$ as a fixed point, so
the selection mechanism is operative for all $h>0$.

Section~\ref{sec:numerics}'s deep-ladder extrapolations along the
joint-limit schedule converge to a limit cleanly distinct from
$P_{\rm FR}$: at $(\tau,\gamma)=(2,0.098)$ the continuum estimate is
$\deficit_\infty=0.268\pm 0.010$, far from $0$. The economic content of
the main text is therefore unaffected by the appendix's degeneracy;
the appendix simply records the technical reason why the limit
construction is the principled route to the PR equilibrium under a
binary payoff.

\section{Reproducibility}

All data, scripts and certified solutions are in
\texttt{projects/REZN/solved\_fixed\_points/lowtau/} (the 100-cell
low-$\tau$ archive) and
\texttt{projects/REZN/solved\_fixed\_points/dd\_k3\_overnight/emin15/}
(the 40-cell mid-$\tau$ archive). Solver code:
\texttt{projects/REZN/solver\_code/highprec\_dd\_qd/}; metrics computation in
\texttt{lowtau\_metrics.py}; report builder in \texttt{lowtau\_report.py}.

\begin{thebibliography}{9}
\bibitem{GS1980} Grossman, S.J. and Stiglitz, J.E. (1980). On the
impossibility of informationally efficient markets. \emph{American
Economic Review} 70(3): 393--408.
\bibitem{Vives2008} Vives, X. (2008). \emph{Information and Learning in
Markets}. Princeton University Press.
\end{thebibliography}

\end{document}
